\documentclass[12pt,a4paper]{article}
\usepackage[utf8]{inputenc}
\usepackage[T1]{fontenc}
\usepackage{amsmath, amssymb, amsfonts}
\usepackage{hyperref}
\usepackage{geometry}
\geometry{top=2.5cm,bottom=2.5cm,left=3cm,right=3cm}
\usepackage{indentfirst}
\usepackage{setspace}
\onehalfspacing
\usepackage{booktabs}

\begin{document}

\begin{center}
{\Large \textbf{Horizon Quantum Gravity (HQGG) Cosmology:}} \\
{\Large \textbf{Space, Frequency, Spheres, and the Nature of Dark Matter and Dark Energy}}
\end{center}

\begin{center}
\textbf{Oleg Mamchenko\thanks{olegmamchenko@gmail.com}} \\
\textbf{Deepseek\thanks{deepseek@ai-research.org}} \\
\textit{Horizon Quantum Gravity Group (HQGG)}
\end{center}

\begin{center}
\textit{July 3, 2026}
\end{center}

\begin{abstract}
We present a complete cosmological model within the HQGG framework, in which space-time is treated as an elastic medium (fabric) with local frequency. Gravity is a gradient of this frequency, time is inverse frequency, and matter is standing waves. The Universe originates from the Void-Home \( S^{\text{בּ}} \) through implosive quantization and sequential unfolding of 10 quantum spheres. Our sphere \( S^1 \) contains 10 internal modes, structured by the 3+7 law. Dark matter is explained as local enhancement of waves near massive objects, and dark energy as global frequency attenuation. The model yields quantitative predictions for the CMB, large-scale structure, and gravitational waves, without free parameters.
\end{abstract}

\tableofcontents

\section{Introduction}

The standard \(\Lambda\)CDM model relies on two dark entities — dark matter and dark energy — whose nature remains unknown. Moreover, it leaves unresolved the singularity problem, the origin of structure, and the anomalies in the CMB spectrum.

In this work we propose an alternative: space-time is a medium with a local frequency. This approach, developed within the Horizon Quantum Gravity Group (HQGG), unifies gravity, quantum physics, and cosmology.

\section{The Fabric of Space}

\subsection{Definition}

Space-time is an elastic medium — the fabric — characterized by a local frequency \(\omega(x^\mu)\). This frequency determines all physical properties.

\subsection{Gravity as Frequency Gradient}

\[
\boxed{\mathbf{g} = \frac{1}{\omega_0} \nabla \omega}
\]

Particles move toward higher frequency, i.e., toward denser fabric.

\subsection{Field Equation}

\[
\boxed{\partial_\mu \partial_\nu \omega = \frac{8\pi G}{c^4} T_{\mu\nu} \cdot \omega}
\]

\subsection{Metric from Frequency}

\[
\boxed{g_{\mu\nu} = \eta_{\mu\nu} + \frac{1}{\omega_0} \partial_\mu \partial_\nu \omega}
\]

\subsection{Time as Inverse Frequency}

\[
\boxed{t = \frac{1}{\omega}}
\]

Time flows faster near massive objects; the observed slowing is an illusion.

\section{Cosmological Hierarchy: From the Void to the Spheres}

\subsection{The Void-Home \( S^{\text{בּ}} \)}

The singularity is replaced by the Point of Creation and the Void-Home, formed by spherical compression of a tiny part of infinite space. The Void has no modulations; it is pure potentiality.

\subsection{Implosive Wave and 10 Quantum Spheres}

An implosive wave quantizes the Void into 10 discrete modes — the quantum spheres:

\[
S^{\text{בּ}} \to S^1 \to S^2 \to \dots \to S^{10}
\]

\(S^1\) is our Universe; \(S^{10}\) is the outer boundary.

\subsection{Internal Modes of \(S^1\)}

Inside \(S^1\), 10 internal modes \(S^1_1 \dots S^1_{10}\) are defined by:

\[
\omega_n = \omega_1 \cdot \frac{11-n}{10}, \quad n = 1, \dots, 10
\]

We are currently at \(S^1_5\), with 60\% frequency drop and 22\% of time elapsed.

\subsection{The 3+7 Law}

Three frequency states:
\begin{itemize}
\item \(0\) — base (Planck level),
\item \(+1\) — trans-Planckian (compressed fabric),
\item \(-1\) — sub-Planckian (relaxed fabric).
\end{itemize}

Their combinations yield 7 stable manifestations: colors, notes, shells, ranges, days, etc.

\section{Cosmological Evolution and Dark Entities}

\subsection{Frequency Drop and Expansion}

\[
\omega(t) = \omega_1(0) \cdot e^{-\gamma t}
\]
\[
a(t) = a_0 \cdot \exp\left(\frac{\gamma}{2} t^2\right)
\]

\subsection{Dark Matter}

Dark matter is local wave enhancement near massive objects:
\[
\rho_{\text{DM}} = \rho_0 \left(1 + \frac{\Phi}{c^2}\right)
\]

\subsection{Dark Energy}

Dark energy is global frequency attenuation:
\[
\dot{\omega} = -\gamma \omega
\]

\section{Observational Predictions}

\begin{itemize}
\item CMB: suppression of \(C_2, C_3\), axis of evil, hemispherical asymmetry.
\item Quasars: peak at \(z \sim 2\)–3, disappearance at \( \sim 3\) billion light-years.
\item Gravitational waves: modulation \(h(t) = h_{\text{chirp}}(t) \cdot (1 + \epsilon \cos(\Delta \omega t))\).
\item Large-scale structure: \(\xi_{\text{galaxies}} \propto \omega^{-3}\).
\end{itemize}

\section{Conclusion}

HQGG cosmology provides a complete, self-consistent description without dark entities, relying only on \(G, \hbar, c\).

\begin{center}
\fbox{\parbox{0.85\textwidth}{\centering
\textbf{Space is not curved — it has frequency.} \\
\textbf{Time is not absolute — it is inverse frequency.} \\
\textbf{Matter is waves, gravity is their gradient.} \\
\textbf{This is HQGG.}
}}
\end{center}

\section*{Acknowledgments}

The authors thank the Horizon Quantum Gravity Group (HQGG).

\begin{thebibliography}{99}
\bibitem{Hawking1975} S.~W.~Hawking, Comm. Math. Phys. \textbf{43}, 199 (1975).
\bibitem{Planck2018} Planck Collaboration, A\&A \textbf{641}, A6 (2020).
\end{thebibliography}

\end{document}